\section{Numerical Analysis of Many Leaks}
\label{sec:results_issues_gt_3}

This section investigates why the proposed formulation and solver, which perform reliably for $n=2$ and $n=3$, fail to consistently identify correct leak parameters when $n \geq 4$. The analysis is based on a systematic simulation study of the Jacobian matrix, solver convergence behaviour, and the effect of scaling and measurement design. The results indicate that the limitation is not primarily algorithmic, but structural and numerical, rooted in severe ill-conditioning of the nonlinear least-squares problem.

\subsection{Method}\label{sec:cond-method}

The analysis is based on the Jacobian matrix $J(x^*) \in \mathbb{R}^{2m \times 2n}$ evaluated at the true parameter vector $x^* = [L_1, \dots, L_n, c_1, \dots, c_n]^T$.
Since $F(x^*) = 0$ by construction (the measurements are generated from the true parameters), the Jacobian at $x^*$ fully characterises the local sensitivity of the residual to parameter perturbations.

The key diagnostic is the \emph{condition number} $\kappa(J) = \sigma_{\max} / \sigma_{\min}$, where $\sigma_{\max}$ and $\sigma_{\min}$ are the largest and smallest singular values of $J$, obtained via its singular value decomposition $J = U \Sigma V^T$.
The condition number bounds the factor by which perturbations (rounding errors, measurement noise) are amplified when solving the linearised system.
A solver working in IEEE~754 double-precision arithmetic ($\varepsilon_{\mathrm{mach}} \approx 1.1 \times 10^{-16}$) can be expected to lose approximately $\log_{10} \kappa$ significant digits, and convergence becomes unreliable when $\kappa$ approaches $1/\varepsilon_{\mathrm{mach}} \approx 10^{16}$.

A high condition number can arise from two distinct causes.
\emph{Scale-induced} ill-conditioning occurs when the Jacobian columns have very different magnitudes; it can be remedied by diagonal scaling (preconditioning).
\emph{Geometric} ill-conditioning occurs when the columns, regardless of their lengths, point in nearly the same directions, i.e.\ the angle between column pairs approaches zero.
To distinguish the two, the \emph{minimum pairwise column angle} is computed: each column of $J$ is normalised to unit length, and the smallest angle between any two normalised columns is reported.
If the column angles are small, no diagonal scaling can improve the condition number, because scaling changes column lengths but not their directions.

\subsubsection{Experimental Setup}

For each number of leaks $n \in \{1, 2, \dots, 8\}$, ten random leak configurations were generated.
In each configuration, positions $z_i$ were drawn uniformly from $[0.05,\, 0.95]$ with a minimum spacing of $0.05$ between consecutive leaks, and coefficients $c_i$ were drawn uniformly from $[0.01,\, 0.2]$.
For each configuration, $m = 2n$ synthetic measurements were generated by simulating the forward model at inlet conditions $(h_0, q_0)$ linearly spaced with $h_0 \in [1.5,\, 20]$ and $q_0 \in [1.5,\, 12]$.
The pipe parameters were $\theta = 0.2$ and $L = 1$.
The Jacobian was evaluated analytically at the true solution using the backward pass (chain-rule differentiation through the forward model), and its SVD was computed using NumPy.
All reported values are averages over the ten configurations, with standard deviations indicating sample-to-sample variability.

To assess the effect of scaling, five strategies were applied to each Jacobian: column scaling (dividing each column by its Euclidean norm), row scaling, combined row and column scaling, and Ruiz iterative equilibration (50~iterations of alternating row/column normalisation).
The condition number of the column-normalised Jacobian serves as the van der Sluis lower bound---the best achievable by any diagonal column scaling.

The effect of measurement diversity was tested for $n = 4$ by varying the inlet condition ranges across four configurations: narrow ($h_0 \in [2, 4]$, $q_0 \in [2, 4]$), medium ($h_0 \in [1, 10]$, $q_0 \in [1.5, 8]$), wide ($h_0 \in [1.5, 20]$, $q_0 \in [1.5, 12]$), and very wide ($h_0 \in [0.5, 40]$, $q_0 \in [0.5, 18]$).
The effect of measurement count was tested for $n \in \{3, 4, 5\}$ with $m$ ranging from $n{+}1$ to $6n$.


\subsection{Results}\label{sec:cond-results}

\subsubsection{Jacobian Conditioning versus Number of Leaks}

Table~\ref{tab:jacobian-cond} presents the condition number, smallest singular value, and minimum column angle of the Jacobian at the true solution, averaged over ten random configurations per~$n$.

\begin{table}[htbp]
\centering
\caption{Jacobian conditioning at the true solution, averaged over 10 random configurations per $n$. Each configuration uses $m = 2n$ measurements with $h_0 \in [1.5, 20]$, $q_0 \in [1.5, 12]$.}
\label{tab:jacobian-cond}
\begin{tabular}{r r c c c c}
\toprule
$n$ & $dim(\mathbf{x})$ & $\overline{\log_{10} \kappa}$ & $\pm\, \sigma$ & avg $\sigma_{\min}$ & avg min.\ angle ($^\circ$) \\
\midrule
1 &  2 & $ 1.31$ & $0.25$ & $7.01 \times 10^{-2}$ & $56.5 \pm 5.1$ \\
2 &  4 & $ 4.48$ & $0.63$ & $1.51 \times 10^{-3}$ & $\phantom{0}3.2 \pm 3.3$ \\
3 &  6 & $ 6.92$ & $0.72$ & $3.59 \times 10^{-6}$ & $\phantom{0}1.9 \pm 1.2$ \\
4 &  8 & $ 9.49$ & $0.81$ & $4.02 \times 10^{-8}$ & $\phantom{0}0.9 \pm 0.7$ \\
5 & 10 & $10.88$ & $0.64$ & $6.15 \times 10^{-10}$ & $\phantom{0}0.9 \pm 0.5$ \\
6 & 12 & $13.35$ & $1.55$ & $1.02 \times 10^{-10}$ & $\phantom{0}0.6 \pm 0.8$ \\
7 & 14 & $15.13$ & $0.98$ & $3.17 \times 10^{-13}$ & $\phantom{0}0.6 \pm 0.4$ \\
8 & 16 & $16.10$ & $0.61$ & $3.27 \times 10^{-14}$ & $\phantom{0}0.4 \pm 0.3$ \\
\bottomrule
\end{tabular}
\end{table}

The condition number grows by approximately two orders of magnitude per additional leak.
At $n = 3$ ($\log_{10} \kappa \approx 7$), the problem is at the limit of what a double-precision solver can handle reliably.
At $n = 4$ ($\log_{10} \kappa \approx 9.5$), approximately 10 of 16 significant digits are consumed by amplification, and the solver stalls.
By $n = 8$ ($\log_{10} \kappa \approx 16$), the condition number reaches $1/\varepsilon_{\mathrm{mach}}$ and the Jacobian is numerically singular.

The minimum column angle decreases simultaneously: from $56.5^\circ$ at $n = 1$ to below $1^\circ$ at $n \geq 4$.
This indicates that the Jacobian columns for adjacent pipe segments become nearly parallel, and the system loses the ability to distinguish perturbations to neighbouring leak parameters from outlet measurements alone.


\subsubsection{Effect of Scaling}

Table~\ref{tab:scaling} compares the condition number (in $\log_{10}$) under five scaling strategies, averaged over ten configurations per $n$.

\begin{table}[htbp]
\centering
\caption{Average $\log_{10} \kappa$ under different scaling strategies. The column-scaled value equals the van der Sluis lower bound.}
\label{tab:scaling}
\begin{tabular}{r c c c c c}
\toprule
$n$ & Unscaled & Column & Row & Row+Col & Ruiz \\
\midrule
1 & $1.31$ & $0.27$ & $1.05$ & $0.62$ & $0.43$ \\
2 & $4.48$ & $4.31$ & $4.00$ & $3.83$ & $3.78$ \\
3 & $6.92$ & $6.67$ & $6.91$ & $6.66$ & $6.65$ \\
4 & $9.49$ & $9.28$ & $9.42$ & $9.20$ & $9.18$ \\
5 & $10.88$ & $10.74$ & $10.83$ & $10.70$ & $10.68$ \\
6 & $13.35$ & $13.27$ & $13.21$ & $13.14$ & $13.13$ \\
7 & $15.13$ & $15.04$ & $15.05$ & $14.99$ & $14.88$ \\
8 & $16.10$ & $16.11$ & $15.98$ & $16.00$ & $15.88$ \\
\bottomrule
\end{tabular}
\end{table}

For $n = 1$, scaling provides a meaningful improvement (from $10^{1.3}$ to $10^{0.3}$), indicating that the $n = 1$ condition number has a scale-related component.
For $n \geq 3$, no strategy reduces $\log_{10} \kappa$ by more than $0.3$, and the column-scaled (geometric) condition number is essentially equal to the unscaled one.
This confirms that the ill-conditioning for $n \geq 3$ is \emph{geometric}: the Jacobian columns are nearly collinear, and scaling cannot change column directions.


\subsubsection{Effect of Measurement Diversity}

Table~\ref{tab:diversity} shows how the range of inlet conditions affects conditioning for $n = 4$ with $m = 8$ measurements.

\begin{table}[htbp]
\centering
\caption{Effect of inlet condition diversity on conditioning ($n = 4$, $m = 8$), averaged over 10 configurations.}
\label{tab:diversity}
\begin{tabular}{l l l c c}
\toprule
Label & $h_0$ range & $q_0$ range & $\overline{\log_{10} \kappa}$ & $\pm\, \sigma$ \\
\midrule
Narrow    & $[2,\, 4]$     & $[2,\, 4]$     & $10.85$ & $0.61$ \\
Medium    & $[1,\, 10]$    & $[1.5,\, 8]$   & $\phantom{0}8.95$ & $1.29$ \\
Wide      & $[1.5,\, 20]$  & $[1.5,\, 12]$  & $\phantom{0}8.88$ & $1.19$ \\
Very wide & $[0.5,\, 40]$  & $[0.5,\, 18]$  & $\phantom{0}8.86$ & $0.92$ \\
\bottomrule
\end{tabular}
\end{table}

Narrow inlet conditions (small variation in $h_0$ and $q_0$) increase the condition number by nearly two orders of magnitude compared to wider ranges.
However, widening the range beyond ``Medium'' gives diminishing returns: the condition number plateaus around $10^{8.9}$ regardless of further diversification.
This plateau reflects the intrinsic ill-conditioning of the problem---once measurements provide sufficient excitation to probe the system, additional diversity cannot resolve the near-parallel Jacobian columns.


\subsubsection{Effect of Number of Measurements}

Table~\ref{tab:more-meas} shows how the condition number changes with the number of measurements $m$ for $n \in \{3, 4, 5\}$.

\begin{table}[htbp]
\centering
\caption{Effect of measurement count $m$ on conditioning, averaged over 10 random configurations. The entry $m = n$ is omitted as the square Jacobian becomes singular for most configurations.}
\label{tab:more-meas}
\begin{tabular}{r r c c}
\toprule
$n$ & $m$ & $\overline{\log_{10} \kappa}$ & $\pm\, \sigma$ \\
\midrule
\multirow{5}{*}{3}
 &  4 & $7.04$ & $0.92$ \\
 &  6 & $6.99$ & $0.94$ \\
 &  9 & $6.61$ & $0.98$ \\
 & 12 & $6.47$ & $0.89$ \\
 & 18 & $6.46$ & $0.86$ \\
\midrule
\multirow{5}{*}{4}
 &  5 & $9.02$ & $1.01$ \\
 &  8 & $8.88$ & $1.19$ \\
 & 12 & $8.50$ & $1.01$ \\
 & 16 & $8.48$ & $0.98$ \\
 & 24 & $8.37$ & $1.03$ \\
\midrule
\multirow{5}{*}{5}
 &  6 & $11.69$ & $0.89$ \\
 & 10 & $10.89$ & $0.65$ \\
 & 15 & $10.32$ & $0.84$ \\
 & 20 & $10.30$ & $0.80$ \\
 & 30 & $10.09$ & $0.84$ \\
\bottomrule
\end{tabular}
\end{table}

For each $n$, increasing the measurement count from $n{+}1$ to $6n$ reduces $\log_{10} \kappa$ by less than one unit.
The largest improvement occurs in the first doubling (from $m \approx n$ to $m \approx 2n$); beyond $m \approx 3n$ the condition number plateaus.
Additional measurements increase $\sigma_{\max}$ (more signal) but cannot lift $\sigma_{\min}$, because the near-null directions of the Jacobian are inherent to the physics: no linear combination of outlet measurements can resolve the compound effect of shifting adjacent leak positions in opposite directions.


\subsection{Discussion}\label{sec:cond-discussion}

The results show that the inverse problem of identifying $n$ leaks from inlet/outlet measurements is fundamentally limited by the conditioning of the Jacobian.
Each additional leak increases $\log_{10} \kappa$ by approximately two units, meaning the condition number grows by roughly two orders of magnitude per leak.
This exponential growth is consistent across all ten random configurations tested per $n$, with moderate standard deviations ($0.6$--$1.6$ in $\log_{10}$ scale), suggesting the trend is robust and not an artefact of a particular leak arrangement.

The root cause is geometric: the Jacobian columns corresponding to adjacent pipe segments become nearly parallel as $n$ grows.
For $n = 4$ the minimum column angle is below $1^\circ$, meaning that the effect of perturbing one segment length is nearly indistinguishable from perturbing its neighbour.
This collinearity arises from the recursive structure of the forward model.
The head and flow at position $i$ depend on all upstream parameters through a chain of $2 \times 2$ transfer matrices $A_k$.
For the outlet measurement, the Jacobian column for parameter $k$ involves the matrix product $A_n A_{n-1} \cdots A_{k+1} B_k$, and adjacent columns differ only by one factor in this product.
As $n$ grows and segments become shorter, adjacent $A_k$ matrices become closer to the identity, making adjacent columns nearly identical.

Scaling strategies cannot address this because they modify column magnitudes, not directions.
More measurements provide diminishing returns because they cannot probe directions that are inherently invisible at the outlet.
Measurement diversity helps to a limited degree by providing varied operating conditions, but the improvement saturates well above the threshold needed for reliable convergence.

These findings impose a practical limit on the identifiability of leaks from inlet/outlet data.
For $n \leq 3$, the condition number ($\sim\! 10^{7}$) is at the edge of double-precision feasibility, and solvers converge reliably.
For $n = 4$ ($\kappa \sim 10^{9.5}$), approximately 10 significant digits are lost to ill-conditioning, and convergence becomes unreliable.
For $n \geq 5$ ($\kappa > 10^{10}$), the problem is effectively unsolvable in standard arithmetic with inlet/outlet measurements alone.

The primary limitation of this analysis is that a single pipe configuration ($\theta = 0.2$, $L = 1$) was used.
Different friction parameters or pipe lengths would change the absolute magnitudes of the Jacobian entries.
However, the recursive matrix-product structure that drives the exponential growth is independent of these parameters, so the qualitative scaling is expected to persist.
Additionally, the evenly spaced leak configurations used here represent a favourable case (maximum separation between adjacent leaks); clustered configurations would likely exhibit even worse conditioning.
